%
%============= CHAPTER BREAK =================%
%
\chapter{Apollonius Dyscolus (ca 138--160 CE)}%
\vspace*{3mm} %
\justifying %

\noindent{}Precise dates are elusive for the famous Alexandrian grammarian, Apollonius Dyscolus, %
though he is known to have stayed briefly in Rome.\endnote{%
Montanari \& Montana 2022:1082.} %
His son, Aelius Herodian, a “worthy successor” who often borrowed from and even cited his father,\endnote{%
Montanari \& Montana 2022:1083.} %
lived in Rome for an extended time. %
Herodian worked under the patronage of Marcus Aurelius %
and dedicated his monumental \textit{General Prosody} to the emperor.\endnote{%
Montanari \& Montana 2022:1083.} %
An Aurelian range for the son, and Antonine range for the father, %
are thus reasonable even if somewhat imprecise ranges. %
Both wrote voluminously and were often cited in their own generations and in later centuries %
for their extensive work on the technical dimensions of Greek language. %

The following analysis hinges on two works by Dyscolus, \textit{On Adverbs} / \textit{De adverbiis} / Περὶ ἐπιρρημάτων, %
and \textit{On Syntax} / \textit{De constructione} / Περὶ σύνταξεως, %
which are largely intact and generally considered authentic. %
Herodianʼs authorship is another matter entirely. %
For a fuller discussion of the authenticity and inauthenticity of works %
attributed to Herodian, see the opening to the later chapter on him. %
In this chapter, we cover only words that are clearly and first attested by his father, %
supplementing this analysis with references to Herodianʼs works that may be authentic or dubious. %
Quotations from works of dubious attribution are indicated with “[Ps-]Herodian” as the author. \\[2mm] %

\noindent{}\textbf{βρύχω} / \textit{brucho} / “to gobble|devour|gnaw” is a postclassical %
Greek form of the classical Greek βρύκω / \textit{brukō}, a lemma commonplace from Hesiod onward (736 TLG hits), %
regularly used by non-Christian and Christian authors alike. %
The variant spelling with a \textit{chi} in place of the \textit{kappa} is later and rarer (21 TLG hits). %
This particular lemma first appears in Dyscolusʼ treatise \textit{On Adverbs}.\endnote{%
\textit{GG} 2.1:197.%
} 

\vspace*{3mm} 
 \begin{quote}καὶ παρὰ ῥῆμα σχηματίζεται, %
 ἐκ μὲν βαρυτόνου τοῦ \textbf{βρύχω} τὸ βρυχηδόν· \\[1mm] %
 [Words ending in -\textit{don}] are also formed from verbs; %
 from the uncircumflexed verb “\textbf{to gnaw}” comes “gnawingly”. \end{quote} % 
 \vspace*{3mm}

A generation later, his son uses the same, unusual spelling in his \textit{General Prosody},\endnote{%
\textit{GG} 3.1:436--437. % %
} %
the monumental work dedicated to Marcus Aurelius. %

\vspace*{3mm} 
 \begin{quote}Τὰ εἰς χω δισύλλαβα ἀρχόμενα ἀπὸ συμφώνου μὴ διὰ τοῦ η παραληγόμενα φύσει μακρᾷ βαρύνεται· 
 εἰ δέ τι περιεσπάσθη, τοῦτο ἐκ βαρυτόνου παρήχθη ὡς ἐπὶ τοῦ τρέχω τροχῶ ... %
 βαρύνεται δὲ ταῦτα, στείχω, τεύχω, ψύχω, τρύχω, σμύχω, \textbf{βρύχω} \\[1mm] %
 Disyllabic words beginning with consonants, ending in χω, %
 when not lengthened by η, are naturally unaccented on the long [syllable]. %
 But if a certain one is circumflexed, it was derived from an unaccented word, %
 as with \textit{trecho} for \textit{trochō} ... %
 These are unaccented: \textit{steicho}, \textit{teucho}, %
 \textit{psucho}, \textit{trucho}, \textit{smucho}, \textbf{to gnaw} / \textit{\textbf{brucho}} \end{quote} % 
 \vspace*{3mm} 

\noindent{}The lemma also appears in the \textit{Partitiones} of [Ps-]Herodian,\endnote{Boissonade 1819:11.} %
in a list of words starting with the letters βρυχ / \textit{bruch}: %

\vspace*{3mm} 
 \begin{quote}\textbf{βρύχω}, ἐπὶ ἀνθρώπου· βρυχῶμαι, ἐπὶ λέοντος· καὶ βρύχημα. \\[1mm]
 \textbf{To gnaw}, [said] about a human; to roar, [said] about a lion; and roar. %
 \end{quote} % 
 \vspace*{3mm}

Despite the Alexandria-based LXX using the classical and more common spelling five times,%
\endnote{Ps 34:16, 36:12, 111:10, Job 16:9, Lam 2:16.} %
the only occurrence of the lemma in the NT follows the rare spelling %
used by two blood-related, leading grammarians of the mid-2nd century. %
Even more fascinating in context is the characterization of Stephenʼs would-be killers %
as carnivorous animals with hearts made of saw-like teeth (Acts 7:54): %

\vspace*{3mm} 
 \begin{quote}Ἀκούοντες δὲ ταῦτα διεπρίοντο ταῖς καρδίαις αὐτῶν καὶ \textbf{ἔβρυχον} τοὺς ὀδόντας ἐπ᾽ αὐτόν. \\[1mm]
 Now hearing these things, they were sawing asunder in their hearts and \textbf{were gnawing} the teeth over him. \end{quote} % 
\vspace*{3mm}

The confluence of a new lemma and astonishingly similar surrounding context %
favor not merely the general neologistic influence of Dyscolus, %
but indeed the more immediate influence of a word list akin to that attributed to Herodian. %
This connection also highlights by contrast how English %
translations of διεπρίοντο / \textit{dieprionto} as a bland passive %
or mere emotional state %
are erroneous morphologically (the verb is likely middle) %
and/or oblivious to the animal imagery. Compare: %

\begin{spacing}{0.8}
\begin{small}
\begin{itemize}
\item “they were cut to the heart” (KJV)
\item “they were cut to the quick” (NAS)
\item “they became enraged” (NRS)
\item “they were furious” (NIV)
\end{itemize}
\end{small}
\end{spacing} \\[1.5mm]

\noindent{}For lexical context, the LSJ entry for διαπρίω / \textit{diapriō} opens: %
“-ιοῦμαι, to saw quite through, saw asunder”. %
To wit, Stephenʼs executioners most closely resemble the lexicographic lions of [Ps-]Herodian. \\[2mm] %

\noindent{}\textbf{γαμίζω} / \textit{gamidzō} / “to marry” | “to marry off” (78 TLG hits) might surprise some readers as a neologism candidate, %
given that it occurs seven times in canonical NT texts: Mt 22:30, 24:38, Mk 12:25, Lk 17:27, 20:35, 1Co 7:38 (\textit{bis}). %
But other lemmas with similar meaning appear earlier in Greek literature, %
including the verb γαμίσκω / \textit{gamiskō} / “to marry” | “to marry off” (60 TLG hits from Aristotle on), %
and the more ancient and far more prevalent cognate γαμέω / \textit{gameō} / %
“to marry” | “to marry off” (8000+ TLG hits from Homer on). %
The first usage of γαμίζω / \textit{gamidzō} outside early %
Christian literature is in Dyscolusʼ treatise \textit{On Syntax} / \textit{De constructione}. %
In the literary context, %
he explains the difference between verbs ending in a simple \textit{omega} and verbs %
ending in \textit{iota-zeta-omega}, %
noting the latter convey the sense of giving something to someone else. %

\vspace*{3mm} 
 \begin{quote}Τὴν αὐτὴν ἔχει διαφορὰν καὶ τὸ γαμῶ πρὸς τὸ \textbf{γαμίζω}· %
 ἔστι γὰρ τὸ μὲν πρότερον γάμου μεταλαμβάνω, τὸ δὲ %
 \textbf{γαμίζω} γάμου τινὶ μεταδίδωμι· τό γε μὴν γαμῶ παθητικῶς κλίνεται, %
 καθὸ τὰ τῆς διαθέσεως πρὸς σύνταξιν ἐμψύχου προσώπου κατεγίνετο. (\textit{Constr}. 153)\endnote{%
 \textit{GG} 2.2:400. Note that this sentence is placed in parentheses in the critical edition, %
 and that manuscript B omits the second occurrence of the lemma and the preceding four words, %
 i.e., γάμου μεταλαμβάνω τὸ δὲ γαμίζω.} \\[1mm]
 The same distinction holds also for “to marry” relative to “\textbf{to marry off}”. %
 For the former is to partake of marriage, %
 but “\textbf{to marry off}” is to give over in marriage to someone. %
 Even so, “to marry” can be inflected as passive, %
 just as matters of function are constructed following the syntax of an animate character. %
 \end{quote} % 
 \vspace*{3mm}

Though the lemma occurs in canonical Luke 17:27 and 20:35, %
and was reconstructed for Marcionʼs \textit{Evangelion} by some editors, %
we conclude it was entirely missing from the \textit{Evangelion}.\endnote{%
Unattested by any witnesses, yet restored for *Ev 17:27 by \textit{HZKN}, but not \textit{VTsB}, %
and possibly attested and restored at *Ev 20:35 by \textit{ZVBRKN}, but not by \textit{H}. 
As with *Ev 20:35, *Ev 20:34 hinges on Tertullianʼs attestation, %
both passive constructions: “The sons of this age marry and are married” / \textit{filii huius aevi nubunt et
nubuntur} (\textit{Marc}. 4.38.8); “neither marry nor are married” / \textit{neque nubere neque nubi} (\textit{Marc}. 4.38.5).
Neither attestation establishes the lemma γαμίζω, contrary to \textit{CPCNT}, which claims it is attested for *Ev 20:35.
} %
It is also notable that Vaticanus and ms 124 attest the older and more commonplace lemma (here γαμίσκονται / \textit{gamiskontai} not %
γαμίζονται / \textit{gamidzontai}) for Luke 20:35. 
Likewise, Marcion's \textit{Apostolos} only has γαμέω / %
\textit{gameō} (*1Cor 7:39),\endnote{%
Missing from \textit{CPCNT}.} %
while γαμίζω / \textit{gamidzō} occurs twice in the immediately preceding canonical verse (1Cor 7:38). %

A pattern begins to emerge. Marcionʼs \textit{Evangelion} and \textit{Apostolos} %
speak of marriage in more ascetic and gender egalitarian terms, %
yet use the most traditional and commonplace word form to do so. %
The canonical NT Gospels and Paulines re-institute patriarchal marriage conventions, %
and yet make use of what appears to be a new word form introduced and/or popularized by Dyscolus. %
This suggests that the final form of the marriage teachings associated with both Jesus and Paul %
were revised and supplemented in a post-Hadrianic setting. \\[1.5mm] %

\noindent{}\textbf{πεισμονή} / \textit{peismonē} / “persuasion|confidence” appears first in a non-Christian %
text in Dyscolusʼ treatise \textit{On Syntax} / \textit{De constructione}. %

\vspace*{2mm} 
 \begin{quote}Οὕτως δ’ οἶμαι ἔχειν καὶ τὸ πείθομαι ὑπὸ σοῦ, %
 καὶ οὐκ ἄλλο τι σημαίνει τὸ πείθομαι σοί %
 ἢ τὴν ἐξ ἀλλήλων πρὸς ἀλλήλους γινομένην \textbf{πεισμονήν}. (\textit{Constr}. 186)\endnote{\textit{GG} 2.2:429.}\\[1mm] %
 So I think it is also [like this with] the [expression] “I am persuaded by you”, %
 and “I am persuaded by you” means nothing else %
 than the \textbf{persuasion} happening from each other toward each other.\end{quote} % 
 \vspace*{2mm}

In Christian texts outside the NT, %
the lemma appears first in Justin Martyrʼs \textit{Dial.Trypho} (ca 155--167 CE): %

\vspace*{2mm} 
 \begin{quote}Πολλὰς μὲν οὖν καὶ ἑτέρας προφητείας ἔχοντες εἰπεῖν, ἐπαυσάμεθα, %
 αὐτάρκεις καὶ ταύτας εἰς \textbf{πεισμονὴν} τοῖς τὰ ἀκουστικὰ καὶ νοερὰ ὦτα %
 ἔχουσιν εἶναι λογισάμενοι καὶ νοεῖν δύνασθαι αὐτοὺς ἡγούμενοι ὅτι οὐχ %
 ὁμοίως τοῖς μυθοποιηθεῖσι περὶ τῶν νομισθέντων υἱῶν τοῦ Διὸς, %
 καὶ ἡμεῖς μόνον λέγομεν ἀλλ’ οὐκ ἀποδεῖξαι ἔχομεν. (\textit{Dial.Trypho} 53.1) \\[1mm] %
 Thus, though having many other prophecies to speak, we pause, %
 reckoning these being sufficient for \textbf{persuasion} to those who have tempered and intellectual ears, %
 and considering them to be able to understand, %
 because they are not similar to the mythmakings about those held as sons of Zeus, %
 and we alone speak but we do not have to prove. \end{quote} % 
 \clearpage

The lemma appears in the NT just once, %
in canonical Gal 5:8, in a portion evidently missing from its Marcionite counterpart: %

\vspace*{2mm} 
 \begin{quote}Ἐτρέχετε καλῶς· τίς ὑμᾶς ἐνέκοψεν [τῇ] ἀληθείᾳ μὴ πείθεσθαι; %
 ⁸ ἡ \textbf{πεισμονὴ} οὐκ ἐκ τοῦ καλοῦντος ὑμᾶς. (Gal 5:7--8) \\[1mm] %
 You were running virtuously. Who thwarted you not to be persuaded by [the] truth? %
 ⁸ The \textbf{persuasion} is not from the one who calls you. %
 \end{quote} % 
 \vspace*{2mm} 

\noindent{}It also appears in all of the major recensions of the letters of Ignatius of Antioch (Ign3 Ign7 Ign13) %
at \textit{Rom} 3.3,\endnote{Cureton 1849:43.} and around 180 CE, in two fragments of %
Irenaeusʼ \textit{Against Heresies} (4 fr. 17, 5 fr. 15).

The lemma did not show up in searches of major databases of digitized Greek papyri and Greek inscriptions. %
The accrued evidence points to the grammarian Dyscolus as the strongest candidate %
for coining and/or popularizing this neologism, %
one that he or his son Herodian brought to Rome, where Justin also taught as a contemporary. %
It also points to the canonical Pauline redaction to Galatians happening during the reign %
of Antoninus Pius at the earliest, and more likely during or after the time of Marcus Aurelius. %